\documentclass[12pt]{article}
\usepackage[utf8]{inputenc}
\usepackage[T1]{fontenc}
\usepackage{hyperref}
\title{San Miguel Piedras Mixtec Language Resources for K-12 Curriculum Development}
\author{Compiled by Claude Sonnet 5.5 (claude-sonnet-5-5) for Andr\'as Kornai}
\date{}
\begin{document}
\maketitle

\section{Introduction}
San Miguel Piedras Mixtec (ISO 639-3 xtp, Glottocode sanm1294) is an Otomanguean language of Oaxaca, Mexico, described as moribund. Reported speaker numbers conflict (Wikipedia cites 460 speakers in 2000; ASJP lists 240). Wikitongues lists the language as unwritten. Searches found no school programme, primer, dictionary or grammar for it.

\section{Primary Educational Resources}
None found.

\section{Linguistic and Reference Materials}
\begin{itemize}
\item Wikipedia entry (overview, classification, speaker figures): \url{https://en.wikipedia.org/wiki/San_Miguel_Piedras_Mixtec}
\item ASJP word-list entry: \url{https://asjp.clld.org/languages/SAN_MIGUEL_PIEDRAS_MIXTEC}
\item PanLex listing: \url{https://apps.panlex.org/panlinx/ex/2746642}
\end{itemize}

\section{Religious and Community Resources}
Global Recordings Network has a language page for xtp: \url{https://globalrecordings.net/en/language/xtp}. The search result did not show whether recordings exist there. Joshua Project lists the people group: \url{https://joshuaproject.net/people_groups/14709/MX}.

\section{Digital and Audio Resources}
Wikitongues has a language page, which may point to community-contributed video: \url{https://wikitongues.org/languages/xtp/}. Its contents were not verified.

\section{Curriculum Development Context}
Very little published material exists. Any curriculum would need to begin with community-led documentation: recording elderly speakers, agreeing on an orthography, and building word lists. Broader Mixtec resources from other varieties could serve as models, but they are different languages and are not listed here.
\end{document}